% Paquetes básicos
\usepackage{booktabs}
\usepackage{pdfpages}
\usepackage{amsthm}
\usepackage{graphicx}

% Configuración de márgenes y espaciado
\usepackage{geometry}
\geometry{
  a4paper,
  top=2.5cm,
  bottom=2.5cm,
  left=3cm,
  right=3cm
}

% Espaciado de teoremas
\makeatletter
\def\thm@space@setup{%
  \thm@preskip=8pt plus 2pt minus 4pt
  \thm@postskip=\thm@preskip
}
\makeatother

% Configuración de portada
\let\oldmaketitle\maketitle
\AtBeginDocument{\let\maketitle\relax}



% ============================================
% ENTORNOS PERSONALIZADOS PARA BLOQUES (TCOLORBOX)
% ============================================
\usepackage[most]{tcolorbox}

% -- Entorno para ADVERTENCIA --
\newtcolorbox{advertencia}{
  enhanced,
  colback=orange!5!white, % Color de fondo muy sutil
  colframe=orange!75!black, % Color del marco
  fonttitle=\bfseries,
  coltitle=orange!50!black, % Color del título
  arc=0mm, % Esquinas cuadradas
  boxrule=0.5pt, % Grosor del marco completo
  borderline west={2mm}{0pt}{orange!75!black}, % Borde izquierdo grueso
  left=5mm,
  right=5mm,
  top=4mm,
  bottom=4mm,
  boxsep=0mm,
  title={Por tanto, esta es una advertencia:} % El título es fijo
}

% -- Entorno para CONCLUSIÓN --
\newtcolorbox{conclusion}{
  enhanced,
  colback=blue!5!white,
  colframe=blue!75!black,
  fonttitle=\bfseries,
  coltitle=blue!50!black,
  arc=0mm,
  boxrule=0.5pt,
  borderline west={2mm}{0pt}{blue!75!black},
  left=5mm,
  right=5mm,
  top=4mm,
  bottom=4mm,
  boxsep=0mm
}

% ============================================



% Mejoras de hipervínculos
\usepackage{hyperref}
\hypersetup{
  colorlinks=true,
  linkcolor=blue,
  filecolor=magenta,
  urlcolor=cyan,
  pdfpagemode=FullScreen,
}

% Configuración de listados de código
\usepackage{listings}
\lstset{
  basicstyle=\ttfamily\small,
  breaklines=true,
  frame=single,
  numbers=left,
  numberstyle=\tiny,
  captionpos=b
}

% NOTA: Los entornos filus y aisa están definidos en fonts.tex
% NO duplicar definiciones aquí